\section{Theory}\label{sec:theory}

The theory distinguishes three possible roles for party-owned digital participation. First, digital platforms may make parties learning organizations by revealing supporter priorities, local information, and policy proposals at low cost. Second, they may operate mainly as mobilization machines, energizing and legitimating supporters without changing party behavior. Third, parties may respond selectively, learning on some issues or during some phases while filtering or ignoring other forms of voice.

The central hypotheses are that comment issue attention predicts later party attention, that concrete and repeated proposals are more likely to appear in later party outputs, and that responsiveness changes as M5S institutionalizes. The strongest mechanisms and scope conditions remain TBD and should be refined before empirical analysis.
